\subsection{Problem Statement}

\subsubsection{Scenario Description}
Consider an enterprise-level application that processes incoming user requests, such as an Online Banking or E-Commerce system. Before executing any core business operation, each request must traverse multiple sequential processing stages to ensure security, correctness, and system stability. These stages typically include:

\begin{enumerate}
    \item \textbf{Authentication:} Verifying whether the user is properly authenticated.
    \item \textbf{Authorization:} Ensuring that the user has sufficient permissions to access the requested resource.
    \item \textbf{Data Validation:} Sanitizing and validating input data to prevent invalid input and security vulnerabilities (e.g., injection attacks).
    \item \textbf{Logging \& Rate Limiting:} Recording request metadata for monitoring purposes and preventing abuse through request throttling.
\end{enumerate}

Such a processing pipeline is fundamental in real-world systems, where each step represents a distinct responsibility within the overall request lifecycle.

\subsubsection{Traditional Approach and Naïve Design}
A straightforward implementation of this pipeline is to encapsulate all processing steps within a single monolithic class using nested conditional logic (e.g., \texttt{if-else} statements):

\begin{lstlisting}[language=Java, caption={Naïve request processing using nested conditionals}]
public class RequestHandler {
    public void process(Request request) {
        if (authenticate(request)) {
            if (authorize(request)) {
                if (validate(request)) {
                    // Execute actual core logic
                    System.out.println("Executing business logic...");
                }
            }
        }
    }
}
\end{lstlisting}

While this approach may appear simple and effective for small-scale systems, it quickly becomes problematic as the number of processing steps increases.

\subsubsection{Architectural Limitations}
As the system evolves, the monolithic design introduces several critical architectural limitations:

\begin{itemize}
    \item \textbf{Violation of the Open/Closed Principle (OCP):} Introducing new processing steps (e.g., data decryption or auditing) requires direct modification of the \texttt{RequestHandler} class, increasing the risk of unintended side effects.
    
    \item \textbf{Tight Coupling:} The request sender is tightly coupled with concrete processing logic, reducing modularity and limiting reuse of individual components.
    
    \item \textbf{High Complexity and Low Maintainability:} Deeply nested conditional structures significantly increase cyclomatic complexity, making the system harder to understand, test, and debug.
    
    \item \textbf{Rigid Execution Flow:} The processing sequence is hardcoded, making it difficult to dynamically modify, extend, or reorder processing steps at runtime.
\end{itemize}

These limitations highlight the need for a more flexible and modular architectural approach.

\subsubsection{Design Goal}
To address these challenges, the system requires a design mechanism that allows multiple independent handlers to process a request in a sequential and decoupled manner. Such a mechanism should support dynamic composition, extensibility, and maintainability, while eliminating direct dependencies between the request sender and specific processing components.

This naturally leads to the adoption of the Chain of Responsibility pattern as a structured solution to these architectural challenges.